% !TEX program = xelatex
% Замена article.tex в репозитории misis-article-latex.
% Рядом должен лежать misisarticle.cls. Сборка: tectonic article.tex
% или GitHub Actions (windows-latest).

\documentclass{misisarticle}

\grnti{28.23.15}
\asjc{1707}
\oecd{1.02}

\title{Мультимодальная 3D-сегментация объектов произвольной формы в системах измерения весогабаритных характеристик}

\misisauthor{П.~А.~Петрова}{кафедра инженерной кибернетики НИТУ «МИСиС»}{Москва, Россия}{m2306412@edu.misis.ru}
\misisauthor{С.~С.~Мишуров}{кафедра инженерной кибернетики НИТУ «МИСиС»}{Москва, Россия}{mishurov.s@misis.ru}

\setabstract{Работа посвящена решению задачи пространственной сегментации разнотипных объектов (коробок, цилиндров, нестандартной тары) в системах автоматизированного измерения весогабаритных характеристик (ВГХ). Основное внимание уделено повышению точности локализации границ крупногабаритных и сложеннопрофилированных грузов, а также обеспечению робастности алгоритмов к оптическим артефактам и отражающим характеристикам материалов поверхности. В работе предложен и апробирован двухпотоковый мультимодальный метод объединения (fusion) пространственных координат, получаемых с помощью ToF-сенсора, и текстурно-спектральных данных RGB-диапазона на базе глубоких нейросетевых архитектур PointNet++ и RandLA-Net. Оценка эффективности разработанного подхода проведена на базе программно-аппаратного комплекса Garpix 3D Scan Visioner с использованием метрик mIoU и F1-score при обеспечении времени инференса модели в пределах 500 мс.}

\setkeywords{весогабаритные характеристики (ВГХ), облако точек (ОТ), ToF-камера, RGB-камера, сверточная нейронная сеть (CNN), программно-аппаратный комплекс (ПАК)}

\begin{document}
\maketitle

\section{Введение}

Складская обработка всё чаще требует знать весогабаритные характеристики груза сразу при приёмке: длину, ширину, высоту и массу. От этих величин зависят тарификация перевозки, место в системе управления складом и карточка товара. Измерение вручную на таком потоке не удерживает ни темп, ни повторяемость.

Автоматические комплексы обычно исходят из того, что предмет близок к прямоугольному параллелепипеду, и считают описывающий прямоугольный ящик по облаку точек. Для цилиндра, мягкой упаковки и деформированной тары такой ящик систематически завышает габарит. Ошибка появляется ещё раньше, на сегментации: если к маске объекта прилипла платформа или, наоборот, срезан край груза, размеры уже не восстановить.

На практике форма --- не единственный источник сбоя. На платформу могут поставить несколько мест сразу. Грузы перекрывают друг друга. Картон, ткань, зеркальный металл и прозрачный пластик по-разному возвращают инфракрасный импульс времяпролётной камеры: в облаке появляются дыры, ложные отражения и мультипатные помехи. Цвет и фактура при этом часто остаются читаемыми на обычном снимке. Поэтому для границ крупного и сложного груза нужны оба канала, а не только глубина.

Рассматривается программно-аппаратный комплекс Garpix 3D Scan Visioner. Объект стоит на весовой платформе. ToF-камеры строят плотное облако по времени пролёта инфракрасного импульса, RGB-камера снимает цвет. По данным производителя один цикл измерения занимает 1--2~с, предел взвешивания --- до 3000~кг, погрешность габарита --- не более 50~мм; результат уходит в учётные системы по USB, Ethernet или Wi-Fi~\cite{garpix}. Встроенный модуль уже отделяет объект от платформы и посторонних предметов. Для груза произвольной формы этот шаг остаётся главным ограничением точности.

В работе предлагается двухпотоковое объединение координат ToF и цвета RGB. Один поток --- PointNet++~\cite{pointnetpp}: иерархические локальные признаки, по которым проходит граница коробки, цилиндра и нестандартной тары. Второй --- RandLA-Net~\cite{randlanet}: случайное прореживание большого облака, чтобы вывод одной сцены уложился в 500~мс и не съел цикл измерения. Ниже разобраны существующие методы, поставлена задача и описан сам способ слияния.

\section{Анализ работ}

\subsection{Системы измерения весогабаритных характеристик}

Ультразвуковые датчики и лазерные дальномеры дёшевы и быстры, но снимают одну или несколько осей. Полной формы они не дают и рассчитаны на правильную геометрию. Структурированный свет и стереозрение восстанавливают поверхность плотнее, однако засветка, ограниченная база и зеркальные или прозрачные материалы оставляют в облаке пропуски. ToF-камера последнего поколения компактнее лазерного сканера, работает при разной освещённости и сразу измеряет дальность, без триангуляции. Слабые места те же, что у активного оптического канала: край объекта и зеркальная поверхность~\cite{garpix}.

Семейство Garpix 3D Scan совмещает ToF, цвет и весовую платформу. Масса сама по себе границу не проводит, но расхождение массы и видимого объёма --- повод проверить, не лежит ли на платформе несколько мест. Пространственная задача от этого не снимается: прежде чем считать ящик, нужно отделить точки груза от пола, рамы и соседнего предмета.

\subsection{Нейросетевая сегментация облаков точек}

Облако --- неупорядоченное множество, и обычная двумерная свёртка к нему не применяется. PointNet обрабатывает каждую точку общим многослойным перцептроном и собирает сцену симметричным max pooling, поэтому перестановка точек результат не меняет~\cite{pointnet}. Блоки T-Net оценивают поворот облака. Для сегментации глобальный вектор склеивается с признаками точки и подаётся в классификатор. Соседей точка не видит: два соприкасающихся места сеть легко склеивает в одно.

PointNet++ переносит на облако идею растущего рецептивного поля~\cite{pointnetpp}. Опорные точки выбираются как наиболее удалённые (farthest point sampling), вокруг каждой берётся шар фиксированного радиуса или $k$ ближайших соседей, и локальный набор снова проходит через PointNet. Следующий уровень работает уже с этими центрами. Так сеть держит и стык упаковок, и силуэт паллеты и меньше зависит от того, что вблизи ToF-камеры плотность точек выше, чем у дальней грани.

DGCNN на каждом слое заново соединяет ближайших соседей, но уже в пространстве признаков, а не только в метрах~\cite{dgcnn}. Признак ребра задаётся центром и разностью с соседом:
\begin{equation}
e_{ij} = h(x_i,\, x_j - x_i),
\label{eq:edgeconv}
\end{equation}
где $h$ --- обучаемое отображение. Точки с похожей геометрией сходятся, даже если между ними в сцене есть зазор. Для плотно составленных мест это сильнее PointNet++, но перестройка графа дороже иерархической выборки.

Для большого облака узкое место --- выбор опорных точек. RandLA-Net заменяет его случайной выборкой и сводит стоимость прореживания к константе~\cite{randlanet}. Потерянные детали возвращает блок локального агрегирования: кодирование положения точки относительно соседей и пулинг с вниманием. Для непрерывного потока измерений это прямой способ уложиться во время вывода.

Point Transformer считает внимание между точками и кодирует положение многослойным перцептроном от относительных координат~\cite{pointtransformer}. Качество на публичных бенчмарках высокое, потому что точка видит не только окрестность. Квадратичная цена внимания не оставляет запаса при плотном скане и бюджете 500~мс на встроенном графическом процессоре.

Другой путь к скорости --- развернуть облако в картинку. SqueezeSeg и RangeNet++ строят карту дальности и сегментируют её двумерной свёрточной сетью~\cite{squeezeseg,rangenet}. Вывод быстрый, но две точки на одном луче попадают в один пиксель. RangeNet++ возвращает метки в исходное облако поиском соседей по карте дальности. Воксельные сети, напротив, режут объём на сетку и считают разреженные трёхмерные свёртки только по занятым ячейкам. Цвет на такую сетку ложится устойчиво, однако шаг сетки сам ограничивает точность контура.

Отдельный ход --- не учить трёхмерную сеть с нуля, а спроецировать на облако маски готовой двумерной модели. SAM3D так переносит Segment Anything в 3D и снимает потребность в большой ручной разметке облаков~\cite{sam3d}. Каскад тяжёлого двумерного вывода и обратной проекции плохо помещается в цикл Visioner, если не сокращать модель отдельно.

Отсюда критерии отбора для комплекса ВГХ. Вывод одной сцены --- не более 500~мс на типовом встроенном графическом процессоре. Архитектура должна принимать и координаты, и цвет. Нужен перенос весов с крупных публичных облаков, потому что своя размеченная выборка склада мала. Нужна открытая реализация на распространённом каркасе обучения, чтобы модуль можно было встроить в программное обеспечение комплекса.

\subsection{Объединение ToF и RGB}

Раннее слияние дописывает к точке цвет до какой-либо сети. Для этого трёхмерная точка проецируется на снимок. Способ простой и весь дальнейший расчёт идёт по одному шестимерному входу, но ошибка калибровки сразу попадает в признаки. Позднее слияние учит каналы раздельно и соединяет уже ответы. К сдвигу камер оно устойчивее, однако сеть не видит, как цвет и форма связаны на нижних слоях. Промежуточное слияние держит две ветви и меняется картами признаков через конкатенацию или внимание.

MVX-Net ведёт воксельную ветвь по облаку и двумерную свёрточную ветвь по изображению, а соединяет их точечным вниманием~\cite{mvxnet}. PointPainting сначала получает семантические метки двумерной сети и «рисует» ими точки, которые затем читает трёхмерный детектор~\cite{pointpainting}. Оба метода показывают, что цвет полезен именно как признак точки, а не как отдельная картинка на выходе. Кросс-модальное внимание идёт дальше: признаки одного канала становятся запросом к другому. Качество выше, расход памяти тоже, и для бюджета 500~мс это запасной, а не основной вариант.

Синтетические сцены дополняют редкую реальную разметку. Для склада подходят генераторы с явной геометрией, в частности Blender и CARLA: коробки, цилиндры, смятая тара, несколько мест на платформе. Разрыв между синтетикой и живым сканом сокращается адаптацией домена, а не подменой реального облака.

\subsection{Смежные работы}

Гончаров сравнивает классическую геометрическую реконструкцию простых объектов и сквозную нейронную сеть~\cite{goncharov}. Ермоленко разбирает, как методы стереозрения восстанавливают облако по карте диспаритета~\cite{ermolenko}. Оба результата относятся к получению трёхмерной формы. Они не отделяют разнотипный груз от весовой платформы по совместным данным ToF и RGB и не ставят предел времени вывода для измерительного цикла.

Публичные наборы ShapeNet и ScanNet дают предобучение, но не повторяют сцену комплекса: один груз на платформе, промышленный фон, артефакты времяпролётной камеры~\cite{shapenet,scannet}. Нужен метод, который на такой сцене строит маску объекта произвольной формы, использует цвет против оптических сбоев ToF и удерживает вывод в пределах 500~мс. Этому соответствуют PointNet++ и RandLA-Net в двухпотоковой схеме, а не трансформер и не тяжёлая двумерная модель без доработки.

\section{Постановка задачи}

\subsection{Входные модальности и проекция}

На вход одной сцены поступают две модальности комплекса. Облако ToF-сенсора
\begin{equation}
P = \{p_i\}_{i=1}^{N} \subset \mathbb{R}^{3}, \qquad
p_i = (x_i, y_i, z_i)^{\mathsf{T}}
\label{eq:cloud}
\end{equation}
задано в системе координат сенсора. Цветной снимок $I \in \mathbb{R}^{H \times W \times 3}$ приходит с RGB-камеры. Масса с платформы в сегментацию не входит: её используют только как внешнюю проверку того, что на платформе одно грузовое место.

Точка облака сопоставляется пикселю оператором проекции $\pi \colon \mathbb{R}^{3} \to \mathbb{R}^{2}$. При матрице внутренних параметров $K \in \mathbb{R}^{3 \times 3}$ и внешнем ориентировании $[R \mid t]$
\begin{equation}
\begin{pmatrix} u_i \\ v_i \\ 1 \end{pmatrix}
\sim
K \,[R \mid t]\,
\begin{pmatrix} x_i \\ y_i \\ z_i \\ 1 \end{pmatrix}.
\label{eq:proj}
\end{equation}
Цвет точки $c_i = I(u_i, v_i) \in \mathbb{R}^{3}$. Обогащённое облако
\begin{equation}
P_m = \{(p_i, c_i)\}_{i=1}^{N} \subset \mathbb{R}^{6}
\label{eq:pm}
\end{equation}
есть вход сегментации. Точка вне кадра не получает цвет ближайшего пикселя: её цветовые каналы остаются пустыми.

Искомое отображение $f_{\theta} \colon P_m \to Y$ с весами $\theta$ назначает каждой точке метку. Для семантической сегментации $y_i \in \mathcal{C}$, и набор классов --- измеряемый объект, весовая платформа, фон, посторонний предмет. Для нескольких мест на платформе нужна инстанс-метка: $y_i = 0$ у фона и платформы, а значения $1, \ldots, M$ различают экземпляры. Паноптический вариант совмещает оба ответа: фон и платформа семантические, грузы --- отдельные экземпляры. Для ВГХ это основной режим, потому что склеенные места дают один неверный габарит.

\subsection{Метки и функция потерь}

Обучение минимизирует потерю на выборке $\mathcal{D} = \{(P_m^{(k)}, Y^{(k)})\}_{k=1}^{K}$:
\begin{equation}
\theta^{*}
=
\arg\min_{\theta}
\frac{1}{K}
\sum_{k=1}^{K}
\mathcal{L}\!\left(f_{\theta}(P_m^{(k)}), Y^{(k)}\right).
\label{eq:train}
\end{equation}
Точек фона и платформы много больше, чем точек мелкого груза, поэтому одна перекрёстная энтропия $\mathcal{L}_{\mathrm{CE}}$ закрашивает объект фоном. Берётся сумма
\begin{equation}
\mathcal{L}
=
\alpha \mathcal{L}_{\mathrm{CE}}
+
\beta \mathcal{L}_{\mathrm{Dice}},
\label{eq:loss}
\end{equation}
где веса $\alpha$ и $\beta$ фиксируются до обучения. Потеря Дайса прямо тянет пересечение масок:
\begin{equation}
\mathcal{L}_{\mathrm{Dice}}
=
1 -
\frac{2\sum_{i=1}^{N} y_i \hat{y}_i + \epsilon}
{\sum_{i=1}^{N} y_i + \sum_{i=1}^{N} \hat{y}_i + \epsilon}.
\label{eq:dice}
\end{equation}
Здесь $\hat{y}_i \in [0,1]$ --- предсказанная вероятность целевого класса, $\epsilon > 0$ не даёт знаменателю обратиться в нуль.

\subsection{Ограничения, данные и критерии}

Метод обязан выделять объект независимо от того, коробка это, цилиндр, мягкая тара или предмет произвольной формы. Недосегментация срезает габарит, пересегментация прибавляет к нему платформу. Картон, пластик, металл и ткань дают разные артефакты ToF; цвет как раз нужен там, где глубина врёт. Несколько предметов и взаимное перекрытие нужно хотя бы обнаружить и отдать оператору как ошибку укладки. Если экземпляры разделяются уверенно, каждое место измеряется отдельно.

Вывод ограничен 500~мс на сцену. Цикл комплекса --- 1--2~с, и сегментация занимает только его часть: дальше считаются габарит и обмен с учётной системой.

Своя выборка состоит из реальных сканов Visioner и синтетических сцен той же геометрии, собранных в Blender и CARLA. Предобучение и перенос признаков идут с ShapeNet и ScanNet~\cite{shapenet,scannet}. Публичные наборы не заменяют платформу, раму и типичные сбои ToF, поэтому дообучение на сценах комплекса обязательно.

Качество считается по матрице ошибок, отдельно на класс. Глобальная доля верных точек здесь обманчива: фон велик. Для класса $i$
\begin{equation}
\mathrm{Precision}_i = \frac{TP_i}{TP_i + FP_i},
\qquad
\mathrm{Recall}_i = \frac{TP_i}{TP_i + FN_i}.
\label{eq:pr}
\end{equation}
Их баланс --- $F_1$:
\begin{equation}
F_{1,i}
=
\frac{2 \cdot \mathrm{Precision}_i \cdot \mathrm{Recall}_i}
{\mathrm{Precision}_i + \mathrm{Recall}_i}
=
\frac{2\,TP_i}{2\,TP_i + FP_i + FN_i}.
\label{eq:f1}
\end{equation}
Пересечение над объединением
\begin{equation}
\mathrm{IoU}_i
=
\frac{TP_i}{TP_i + FP_i + FN_i},
\qquad
\mathrm{IoU}_i = \frac{F_{1,i}}{2 - F_{1,i}}.
\label{eq:iou}
\end{equation}
Сводная метрика --- среднее по классам без весов объёма:
\begin{equation}
\mathrm{mIoU}
=
\frac{1}{C}
\sum_{i=1}^{C}
\mathrm{IoU}_i.
\label{eq:miou}
\end{equation}
Именно mIoU и $F_1$ принимаются целевыми: обе штрафуют сдвиг контура и у мелкой посылки, и у паллеты. Время вывода контролируется отдельно и не компенсирует провал по маске.

\section{Описание методов}

\subsection{Калибровка и раннее слияние}

ToF- и RGB-модули Visioner стоят в одном корпусе, поэтому внешнее ориентирование ищется одной калибровкой, а не по каждому измерению. По~\eqref{eq:proj} каждая точка получает цвет. Раннее слияние на этом и заканчивается: оба потока читают уже вектор $(x, y, z, R, G, B)$, а не голое облако. Интенсивность отражённого импульса, если сенсор её отдаёт, дописывается отдельным каналом и не смешивается с цветом.

Точки без проекции внутрь кадра не заполняются цветом соседнего пикселя. Иначе блик на раме или этикетка соседнего места становятся признаком чужой точки. Кадры двух сенсоров берутся из одной аппаратной синхронизации: у комплекса общий цикл измерения, и отдельная временная сшивка по меткам не требуется.

Перед сетью облако переводится в систему координат платформы. T-Net внутри PointNet умеет компенсировать поворот, но для измерительной сцены надёжнее явная опора: плоскость платформы и вертикаль. Тогда высота груза не зависит от того, как паллета стоит в зоне измерения.

\subsection{Поток PointNet++}

Первый поток отвечает за границу. PointNet++ строится как иерархия множеств~\cite{pointnetpp}. На уровне $l$ опорные точки выбираются farthest point sampling, чтобы покрытие не схлопнулось в ближнюю к камере грань, где точек больше. Вокруг опоры собираются соседи в шаре, радиус которого растёт с уровнем: нижний шар ловит кромку коробки и ребро цилиндра, верхний --- весь силуэт.

Каждая окрестность проходит мини-PointNet: общий перцептрон по точкам окрестности и max pooling. Во внутреннем представлении точка несёт и координаты, и цвет, поэтому пулинг видит и скачок глубины, и скачок фактуры. Обратный ход интерполяцией возвращает признаки на исходные точки. Без него маска жила бы на прореженном наборе и снова огрубляла габарит.

Этот поток дороже второго. Farthest point sampling квадратичен по числу точек, и на полном скане крупного груза он один способен выйти за 500~мс. Поэтому PointNet++ не обязан читать каждую точку сырого облака: на вход ему подаётся облако после того же случайного прореживания, которым пользуется второй поток, с сохранением соответствия признаков исходным точкам.

\subsection{Поток RandLA-Net}

Второй поток отвечает за время. RandLA-Net не ищет удалённые опоры~\cite{randlanet}. На каждом уровне точки выбираются случайно, стоимость выбора не растёт с квадратом размера облака. Чтобы прореживание не съело кромку, блок локального агрегирования делает две вещи. Кодирование положения складывает относительные координаты соседей с признаками центра: сеть видит, что точка лежит на краю, а не только то, какого она цвета. Пулинг с вниманием усиливает тех соседей, которые согласованы с центром, и глушит выброс ToF --- одиночную точку от мультипатного отражения.

Декодер поднимает признаки обратно к более плотному набору. На выходе у каждой сохранённой точки есть свой вектор, как и у PointNet++. Общая случайная выборка на входе нужна, чтобы оба вектора относились к одним и тем же точкам и их можно было склеить без второй интерполяции.

Случайность выборки при выводе фиксируется. Иначе повторный прогон той же сцены шевелит контур на величину, сравнимую с заявленной погрешностью габарита 50~мм. Для измерения это хуже, чем постоянный небольшой сдвиг.

\subsection{Промежуточное слияние и ответ сети}

Раннее слияние уже смешало цвет и координаты внутри каждого потока. Промежуточное слияние соединяет потоки между собой. Пусть $h^{\mathrm{pp}}_i$ --- признак точки из PointNet++, $h^{\mathrm{rl}}_i$ --- признак той же точки из RandLA-Net. Вектор точки перед классификатором
\begin{equation}
h_i = \bigl[h^{\mathrm{pp}}_i \,\|\, h^{\mathrm{rl}}_i\bigr].
\label{eq:fuse}
\end{equation}
Дальше --- общий многослойный перцептрон и softmax по семантическим классам. Градиент потери~\eqref{eq:loss} идёт в обе ветви, так что поток границ учится не дублировать быстрый поток, а поправлять его там, где контур тонкий.

Инстанс-метки строятся по тем же признакам $h_i$ для точек, которые семантическая голова отнесла к грузу. Точки платформы и фона в экземпляры не входят. Если уверенных экземпляров больше одного или маски пересекаются так, что общий ящик заведомо покрывает два места, сцена помечается как ошибка укладки. Это соответствует правилу комплекса: несколько предметов либо разделяются, либо измерение не выпускается.

Оба потока остаются в одном графе вывода на PyTorch; облако читается средствами Open3D. Отдельной двумерной сегментационной сети, как в PointPainting, в контуре нет: цвет входит ранним каналом, а не готовой маской. Так не появляется второй тяжёлый вывод, который сам по себе не укладывается в 500~мс.

Порядок расчёта одной сцены такой.

\begin{enumerate}
\item Взять синхронную пару «облако ToF --- кадр RGB» и показание весов оставить за пределами сети.
\item Спроецировать точки по~\eqref{eq:proj}, собрать $P_m$ и отбросить цвет у точек вне кадра.
\item Случайно проредить облако фиксированным образом и подать его в PointNet++ и в RandLA-Net.
\item Склеить поточечные признаки по~\eqref{eq:fuse}, получить семантические метки и, если группы точек груза не связаны между собой, инстанс-метки.
\item Если маска объекта одна и вывод занял не больше 500~мс, передать точки объекта модулю габарита. Иначе вернуть оператору отказ.
\end{enumerate}

Целевые показатели этого контура --- mIoU и $F_1$ по классам постановки и время вывода. Сравнение имеет смысл только между вариантами одной и той же сцены Visioner: полный двухпотоковый контур, каждый поток по отдельности и тот же поток без цветовых каналов. Последний вариант показывает, даёт ли RGB устойчивость на зеркальных и прозрачных поверхностях или только увеличивает вход. Числовые значения метрик определяются этим сравнением и в описание метода не закладываются заранее.

\begin{thebibliography}{99}

\bibitem{pointnet} Qi~C.R., Su~H., Mo~K., Guibas~L.J. PointNet: Deep Learning on Point Sets for 3D Classification and Segmentation~// Proceedings of the IEEE Conference on Computer Vision and Pattern Recognition (CVPR).~--- 2017.~--- P.~652--660.

\bibitem{pointnetpp} Qi~C.R., Yi~L., Su~H., Guibas~L.J. PointNet++: Deep Hierarchical Feature Learning on Point Sets in a Metric Space~// Advances in Neural Information Processing Systems (NeurIPS).~--- 2017.~--- Vol.~30.

\bibitem{dgcnn} Wang~Y., Sun~Y., Liu~Z., Sarma~S.E., Bronstein~M.M., Solomon~J.M. Dynamic Graph CNN for Learning on Point Clouds~// ACM Transactions on Graphics.~--- 2019.~--- Vol.~38, No.~5.~--- P.~1--12.

\bibitem{pointtransformer} Zhao~H., Jiang~L., Jia~J., Torr~P., Koltun~V. Point Transformer~// Proceedings of the IEEE/CVF International Conference on Computer Vision (ICCV).~--- 2021.~--- P.~16259--16268.

\bibitem{randlanet} Hu~Q., Yang~B., Xie~L., Rosa~S., Guo~Y., Wang~Z., Trigoni~N., Markham~A. RandLA-Net: Efficient Semantic Segmentation of Large-Scale Point Clouds~// Proceedings of the IEEE/CVF Conference on Computer Vision and Pattern Recognition (CVPR).~--- 2020.~--- P.~11108--11117.

\bibitem{squeezeseg} Wu~B., Wan~A., Yue~X., Keutzer~K. SqueezeSeg: Convolutional Neural Nets with Recurrent CRF for Real-Time Road-Object Segmentation from 3D LiDAR Point Cloud~// Proceedings of the IEEE International Conference on Robotics and Automation (ICRA).~--- 2018.~--- P.~1887--1893.

\bibitem{rangenet} Milioto~A., Vizzo~I., Behley~J., Stachniss~C. RangeNet++: Fast and Accurate LiDAR Semantic Segmentation~// Proceedings of the IEEE/RSJ International Conference on Intelligent Robots and Systems (IROS).~--- 2019.~--- P.~4213--4220.

\bibitem{sam3d} Yang~Y., Wu~X., He~T., Zhao~H., Liu~X. SAM3D: Segment Anything in 3D Scenes~// arXiv preprint.~--- 2023.~--- arXiv:2306.03908.

\bibitem{shapenet} Chang~A.X., Funkhouser~T., Guibas~L., Hanrahan~P., Huang~Q., Li~Z. et~al. ShapeNet: An Information-Rich 3D Model Repository~// arXiv preprint.~--- 2015.~--- arXiv:1512.03012.

\bibitem{scannet} Dai~A., Chang~A.X., Savva~M., Halber~M., Funkhouser~T., Nie{\ss}ner~M. ScanNet: Richly-Annotated 3D Reconstructions of Indoor Scenes~// Proceedings of the IEEE Conference on Computer Vision and Pattern Recognition (CVPR).~--- 2017.~--- P.~5828--5839.

\bibitem{mvxnet} Sindagi~V.A., Zhou~Y., Tuzel~O. MVX-Net: Multimodal VoxelNet for 3D Object Detection~// Proceedings of the IEEE International Conference on Robotics and Automation (ICRA).~--- 2019.~--- P.~7276--7282.

\bibitem{pointpainting} Vora~S., Lang~A.H., Helou~B., Beijbom~O. PointPainting: Sequential Fusion for 3D Object Detection~// Proceedings of the IEEE/CVF Conference on Computer Vision and Pattern Recognition (CVPR).~--- 2020.~--- P.~4604--4612.

\bibitem{garpix} Официальный сайт компании GARPIX: система измерения ВГХ Garpix 3D Scan Visioner. URL: \url{https://garpixscan.com} (дата обращения: 10.05.2025).

\bibitem{goncharov} Гончаров~А.Д. 3D-реконструкция простых объектов. Сравнение традиционного геометрического метода и end-to-end нейронной сети. URL: \url{https://sadekov.su/Articles/kik_2026_1.pdf}.

\bibitem{ermolenko} Ермоленко~И.Д. Анализ эффективности методов стереозрения для восстановления 3D-облака точек по карте диспаритета. URL: \url{https://sadekov.su/Articles/kik_2026_1.pdf}.

\end{thebibliography}

\end{document}
